\documentclass{exam}
\usepackage{../../shah350}

\hypersetup{colorlinks=true, linktoc=section, linkcolor=blue}

\pagestyle{headandfoot}
\firstpageheadrule
\runningheadrule
\firstpageheader{Prof. Lepowsky \\ Linear Algebra}{Homework \#9}{\href{https://github.com/jeevanshah07/MATH350/blob/main/homework/homework9/main.tex}{\underline{Source}} \\ Jeevan Shah}
\runningheader{Linear Algebra \\ Homework \#9}{}{Shah}
\firstpagefooter{}{\thepage}{}
\runningfooter{ }{\thepage}{ }

\printanswers

% Sec. 5.4 #2(b),
% 3 (a), (b), (c) - all proofs
% 6(a),(b),
% 9 (a) (b) - (i), (ii),
% 10 (a) (b),
% 16 (a) (b) - all proofs,
% 18 (a) (b) (c) - all proofs,
% 21 - proof

\begin{document}
  \begin{problem}{5.4.2b}
    Consider $f(x) = x^2 \in W$ so $T(f(x)) = xf(x) = x(x^2) = x^3 \not\in W$. Hence $T(W) \not\subseteq W$ so $W$ is not $T$ invariant.
  \end{problem}

  \begin{problem}{5.4.3}
    \begin{parts}
      \part \begin{proof} 
        \begin{enumerate}[label=\roman*).]
          \item Clearly $\braces{\vec{0}}$ is T-invariant as the only vector in $\braces{\vec{0}}$ is $\vec{0}$ and $T(\vec{0}) = \vec{0}$ for any linear operator. Hence $T\left(\braces{\vec{0}}\right) \subseteq \braces{\vec{0}}$.
          \item $V$ is trivially T-invariant as $T$ has co-domain $V$ since it is a linear operator and hence for all $v \in V,\, T(v) \in V$ by definition and thus $T(V) \subseteq V$.
        \end{enumerate}
      \end{proof}

      \part \begin{proof}
        \begin{enumerate}[label=\roman*.)]
          \item Recall that $N(T) = \braces{v \in V \mid T(v) = \vec{0}}$ and notice $\vec{0} \in N(T)$ since $T(\vec{0}) = \vec{0}$. It follows that for all $x \in N(T)$, $T(x) = \vec{0}$ and so $T(x) \in N(T)$. Hence $T(N(T)) \subseteq N(T)$.
          \item Recall that $R(T) = \braces{T(v) \mid v \in V}$ is the image of $T$. Hence for all $v \in R(T)$, $T(v) = T(T(x))$ for some $x \in V$. However, since $T$ is a linear operator, $v = T(x) \in V$ so $T(v) \in R(T)$ and thus $T(R(T)) \subseteq R(T)$.
        \end{enumerate}
      \end{proof}

      \part \begin{proof}
        Let $\lambda \in \FF$ be an eigenvalue of $T$ and $E_\lambda$ be its corresponding eigenspace. Recall that $E_\lambda = \braces{x \in V \mid T(x) = \lambda x}$. Then for all $x \in E_\lambda$, $T(x) = \lambda x$ by definition, and $\lambda x \in E_\lambda$ since $E_\lambda$ is a subspace. Hence $T(E_\lambda) \subseteq E_\lambda$.
      \end{proof}
    \end{parts}
  \end{problem}

  \begin{problem}{5.4.6}
    \begin{parts}
      \part Notice that if $z = e_1 = (1, 0, 0, 0)$ then 
      \begin{align*}
        T(e_1) &= (1, 0, 1, 1) \\
        T^2(e_1) &= (1, -1, 2, 2) \\
        T^3(e_1) &= (0, -3, 3, 3) \\
        &= 0 \cdot e_1 - 3 \cdot T(e_1) + 3 \cdot T^2(e_1)
      \end{align*}
      Hence $T^3(e_1)$ is a linear combination of $e_1, T(e_1), T^2(e_1)$, and these three vectors are linearly independent (each is outside the span of the previous ones), and thus $\beta = \braces{e_1, T(e_1), T^2(e_1)}$ is an ordered basis for the T-cyclic subspace generated by $z$.

      \part Notice that if $z = x^3$ then 
      \[
        T(z) = 6x \mand T^2(z) = \vec{0}
      \]
      and $T^2(z) = \vec{0} = 0 \cdot z + 0 \cdot T(z)$ is a linear combination of $z, T(z)$. Since $x^3$ and $6x$ are not scalar multiples of each other, $\braces{z, T(z)}$ is linearly independent and hence $\beta = \braces{z, T(z)} = \braces{x^3, 6x}$ is an ordered basis for the T-cyclic subspace generated by $z$.
    \end{parts}
  \end{problem}

  \newpage 

  \begin{problem}{5.4.9}
    \begin{parts}
      \part \begin{enumerate}[label=\roman*.)]
        \item Since $\beta = \braces{e_1, T(e_1), T^2(e_1)}$ is a basis for $W$ and $T^3(e_1) - 3T^2(e_1) + 3T(e_1) + 0 \cdot e_1 = 0$ (so $a_0 = 0$, $a_1 = 3$, $a_2 = -3$), it follows by Thm~5.21(b) that the characteristic polynomial is 
        \[
          f(t) = (-1)^3(3t - 3t^2 + t^3) = -t^3 + 3t^2 - 3t
        \]
        \item We have 
        \[
          [T_W]_\beta = \begin{pmatrix}
            0 & 0 & 0 \\
            1 & 0 & -3 \\
            0 & 1 & 3
          \end{pmatrix}
          \Rightarrow \det([T_W]_\beta - tI) = \begin{vmatrix}
            -t & 0 & 0 \\
            1 & -t & -3 \\
            0 & 1 & 3 - t
          \end{vmatrix}
          = -t(-t(3-t) + 3) = -t^3 + 3t^2 -3t
        \]
      \end{enumerate}

      \part \begin{enumerate}[label=\roman*.)]
        \item Since $\beta = \braces{z, T(z)}$ is a basis for $W$ and 
        \[
          T^2(z) = 0 \cdot T(z) + 0 \cdot z
        \]
        the characteristic polynomial is 
        \[
          f(t) = (-1)^2 t^2 = t^2
        \]
        \item We have that 
        \[
          [T_W]_\beta = \begin{pmatrix}
            0 & 0 \\
            1 & 0 
          \end{pmatrix}
          \Rightarrow \det([T_W]_\beta - tI) = \begin{vmatrix}
            -t & 0 \\
            1 & -t 
          \end{vmatrix}
          = t^2
        \]
      \end{enumerate}
    \end{parts}
  \end{problem}

  \begin{problem}{5.4.10}
    \begin{parts}
      \part Notice that $T = L_A$ for 
      \[
        A = \begin{pmatrix}
          1 & 1 & 0 & 0 \\
          0 & 1 & -1 & 0 \\
          1 & 0 & 1 & 0 \\
          1 & 0 & 0 & 1
        \end{pmatrix}
      \]
      hence we can find the characteristic polynomial for $T$: 
      \begin{align*}
        f(t) = \det(A - tI) &= \begin{vmatrix}
          1 - t & 1 & 0 & 0 \\
          0 & 1-t & -1 & 0 \\
          1 & 0 & 1-t & 0 \\
          1 & 0 & 0 & 1-t
        \end{vmatrix} \\
        &= (1 - t)\begin{vmatrix}
          1 - t & 1 & 0 \\
          0 & 1 - t & -1 \\
          1 & 0 & 1-t 
        \end{vmatrix} \\
        &= (1 - t)\left(\begin{vmatrix}
          1 & 0 \\
          1 - t & -1
        \end{vmatrix} + (1 - t)\begin{vmatrix}
          1 - t & -1 \\
          0 & 1-t
        \end{vmatrix}\right) \\
        &= (1 - t)(-1 + (1 - t)^3) \\
        &= (1 - t)^4 - (1-t) \\
        &= (1 - t)(-t^3 + 3t^2 - 3t) \\
        &= (1 - t) \cdot g(t)
      \end{align*}
      where $g(t) = -t^3 + 3t^2 - 3t$ is the characteristic polynomial for $T_W$. Hence, clearly $g \mid f$.
      \part Let $\gamma$ be the standard ordered basis for $\mathcal{P}_3(\RR)$ and let $A = [T]_\gamma$, so that 
      \[
        A = \begin{pmatrix}
          0 & 0 & 2 & 0 \\
          0 & 0 & 0 & 6 \\
          0 & 0 & 0 & 0 \\
          0 & 0 & 0 & 0
        \end{pmatrix}
      \]
      so $T$ has the same characteristic polynomial as $A$:
      \[
        f(t) = \det(A - tI) = \begin{vmatrix}
          -t & 0 & 2 & 0 \\
          0 & -t & 0 & 6 \\
          0 & 0 & -t & 0 \\
          0 & 0 & 0 & -t
        \end{vmatrix}
        = t^4 = t^2 \cdot g(t)
      \]
      since $A - tI$ is upper triangular and $g(t) = t^2$ is the characteristic polynomial for $T_W$. Hence, clearly $g \mid f$.
    \end{parts}
  \end{problem}

  \begin{problem}{5.4.16}
    \begin{parts}
      \part \begin{proof}
        Let $f$ be the characteristic polynomial of $T$ and assume that $f$ splits over $\FF$. Let $n = \dim(V)$. Then there exist $\lambda_1, \dots, \lambda_n \in \FF$ (not necessarily distinct) such that 
        \[
          f(t) = (-1)^n (t - \lambda_1)(t - \lambda_2) \cdots (t - \lambda_n)
        \]
        Then, let $g$ be the characteristic polynomial of $T_W$ for $W$ some T-invariant subspace. By Thm~5.20 we know that $g \mid f$, i.e., $f(t) = g(t)h(t)$ for some polynomial $h$. Every irreducible factor of $g$ is then an irreducible factor of $f$, and by unique factorization the irreducible factors of $f$ are all linear. Hence $g$ is a product of linear factors, 
        \[
          g(t) = (-1)^k (t - \lambda_{i_{1}})(t - \lambda_{i_{2}}) \cdots (t - \lambda_{i_{k}})
        \]
        where $k = \dim(W)$ and each $\lambda_{i_{j}} \in \braces{\lambda_{1}, \dots, \lambda_{n}}$. Thus $g$ also splits over $\FF$.
      \end{proof}

      \part \begin{proof}
        Let $W$ be a nontrivial $T$-invariant subspace of $V$ and let $g$ be the characteristic polynomial of $T_W$. Since $W \neq \braces{0}$, $\deg(g) = \dim(W) \geq 1$, and $g$ splits by $(a)$, so $g$ has at least one linear factor $t - \lambda$ and hence a root $\lambda \in \FF$. Then $\lambda$ is an eigenvalue of $T_W$, so there is a nonzero $w \in W$ with $T_W(w) = \lambda w$. Since $T(w) = T_W(w) = \lambda w$, $w$ is an eigenvector of $T$ contained in $W$.
      \end{proof}
    \end{parts}
  \end{problem}

  \begin{problem}{5.4.18}
    \begin{parts}
      \part \begin{proof}
        Assume $A$ is invertible. Equivlanently, $\det(A) \neq 0$. Recall that the characteristic polynomial is $f(t) = \det(A - tI)$ so that $f(0) = \det(A)$. However, $f(0) = a_0$ and so $a_0 = \det(A) \neq 0$. \\ Conversely assume $a_0 \neq 0$. As before consider $f(0) = \det(A)$ so that $\det(A) = f(0) = a_0 \neq 0$ and thus $A$ must be invertible.
      \end{proof}

      \part \begin{proof}
        Assume that $A$ is invertible, so that $a_0 \neq 0$ by $(a)$, and let $f$ be the characteristic polynomial for $A$. Then, by the Cayley-Hamilton Theorem we know that $f(A) = O$ so 
        \begin{align*}
          f(A) &= (-1)^n A^n + a_{n-1}A^{n-1} + \cdots + a_1 A + a_0 I_n = O \\
          &\Rightarrow\, (-1)^n A^{n-1} + a_{n-1}A^{n-2} + \cdots + a_2 A + a_1 I_n + a_0 A^{-1} = O \tag{Right multiply by $A^{-1}$} \\
          &\Rightarrow\, (-1)^n A^{n-1} + a_{n-1}A^{n-2} + \cdots + a_2 A + a_1 I_n = -a_0 A^{-1}  \\
          &\Rightarrow\, -\frac{1}{a_0}\left[(-1)^n A^{n-1} + a_{n-1}A^{n-2} + \cdots + a_2 A + a_1 I_n\right] = A^{-1}  
        \end{align*}
      \end{proof}

      \part Notice that $A$ is upper triangular and hence 
      \begin{align*}
        f(t) = \det(A - tI) &= (1 - t)(2 - t)(-1 - t) \\
        &= (2 - 3t + t^2)(-1 - t) \\
        &= -t^3 + 2t^2 + t -2
      \end{align*}
      Hence $a_0 = -2$, $a_1 = 1$, $a_2 = 2$. Now, we have that 
      \[
        A^2 = \begin{pmatrix}
          1 & 2 & 1 \\
          0 & 2 & 3 \\
          0 & 0 & -1
        \end{pmatrix}
        \begin{pmatrix}
          1 & 2 & 1 \\
          0 & 2 & 3 \\
          0 & 0 & -1
        \end{pmatrix}
        = \begin{pmatrix}
          1 & 6 & 6 \\
          0 & 4 & 3 \\
          0 & 0 & 1
        \end{pmatrix}
      \]
      so 
      \[
        A^{-1} = \frac{1}{2}\left((-1)^3 A^2 + 2A + I_3\right) = \boxed{\frac{1}{2}\begin{pmatrix}
          2 & -2 & -4 \\
          0 & 1 & 3 \\
          0 & 0 & -2
        \end{pmatrix}
        }
      \]
    \end{parts}
  \end{problem}

  \begin{problem}{5.4.21}
    \begin{proof}
      Let $T$ be a linear operator on the two-dimensional space $V$ and for $v \in V$ let $W_v$ be the $T$-cyclic subspace generated by $v$. Since $W_v \subseteq V$, $\dim W_v \leq 2$ for every $v$. Either there exists some $v \in V$ with $\dim W_v = 2$, or $\dim W_v \neq 2$ for every $v \in V$. \\
      \underline{Case 1 (there exists $v \in V$ with $\dim W_v = 2$):} Then $W_v$ is a $2$-dimensional subspace of the $2$-dimensional space $V$, so $W_v = V$ and $V$ is a $T$-cyclic subspace of itself. \\
      \underline{Case 2 ($\dim W_v \neq 2$ for every $v \in V$):} Let $v \in V$ be nonzero. Then $\dim W_v \geq 1$ since $v \in W_v$, so $\dim W_v = 1$ and $\braces{v}$ is a basis for $W_v$. Since $T(v) \in W_v$, we have $T(v) = c_v \cdot v$ for some $c_v \in \FF$. Now, let $\braces{a, b}$ be a basis for $V$. Then $a + b \neq 0$ and, by linearity of $T$, $T(a) + T(b) = T(a + b)$ so
      \[
        c_a \cdot a + c_b \cdot b = c_{a + b}(a + b) \Rightarrow 0 = a(c_a - c_{a + b}) + b(c_b - c_{a + b}).
      \]
      Since $a, b$ are linearly independent it follows that $c_a = c_{a + b} = c_b$. Call this common scalar $c$. Then $T(a) = ca$ and $T(b) = cb$, so for any $u = \alpha a + \beta b \in V$, 
      \[
        T(u) = \alpha T(a) + \beta T(b) = c(\alpha a + \beta b) = cu
      \]
      by linearity of $T$. Hence $T = cI$.
    \end{proof}
  \end{problem}
\end{document}